\section{Related Work}
\label{sec:related}

\textbf{Detection on these benchmarks.} Guo et al.~\cite{Guo2023close} introduce HC3 with a RoBERTa detector at 99.82 F1, establishing both the corpus and the accuracy regime the field has worked in since, though Su et al.~\cite{Su2023plus} show that regime narrows to roughly 91.7\% under semantic-invariant tasks. Zhou~\cite{Zhou2024exploiting} reaches competitive DAIGT V2 accuracy with character n-grams and no transformer, consistent with our findings but reported without the surface-content separation that would explain it, and Kubrusly et al.~\cite{Kubrusly2025detecting} likewise report strong bag-of-words results on HC3 without isolating what those features encode. Ansary~\cite{Ansary2026multirepresentation} and Lamsiyah et al.~\cite{Lamsiyah2025daigt} report near-ceiling DAIGT-family results in-distribution only. Two designs sit closest to ours and differ in purpose rather than ingredients. Annepaka et al.~\cite{Annepaka2026synergizing} mix linguistic features with a transformer for 99.45\% on HC3 and Yadagiri et al.~\cite{Yadagiri2025aigenerated} pair DeBERTa with stylometric features, both combining the two channels to raise a score where we hold them apart to measure one. Socolof et al.~\cite{Socolof2024style} detect with style embeddings alone, a motivation we share, though a learned embedding cannot be shown to exclude content the way a hand-built orthographic vector can. Surveys~\cite{Crothers2023machine} catalogue detector families without asking what a corpus is separable by.

\textbf{Competition solutions on DAIGT V2.} The corpus was assembled for a public competition, and the winning entry~\cite{Biswas2024kaggle} and strong public ones~\cite{Zeyadusf2024catch,Zain2025multi,Neo2024didai} lean on character n-grams and ensembles of classical models rather than larger language models. That is usually read as evidence classical models suffice here. Our decomposition suggests a compatible and more actionable reading, that a corpus whose surface arm alone reaches 0.9214 rewards models good at surface statistics, so a leaderboard is partly a measurement of the corpus.

\textbf{Documented artefacts in these corpora.} Tian et al.~\cite{Tian2023multiscale} identify the HC3 whitespace convention and release a cleaning kit, but do not quantify how much of the corpus's separability surface form carries, nor test whether a detector can represent the cue at all, which is the question our tokenisation control answers and which changes the conclusion for BERT. Baidya et al.~\cite{Baidya2026detecting} document HC3's length confound and control it by length-matching, the same confound our two arms share and that Section~\ref{sec:length} closes. Ardeshirifar~\cite{Ardeshirifar2025comparing} standardises punctuation before comparing handcrafted features to deep models without measuring what that removes, and Park et al.~\cite{Park2024investigating} show HC3-trained detectors rely on prompt-specific shortcuts. Where these identify individual cues, we measure the aggregate.

\textbf{Shortcut learning and dataset artefacts elsewhere.} The pattern is not specific to text detection. Torralba and Efros~\cite{Torralba2011unbiased} showed vision datasets are identifiable from their own images, and in natural language inference a model reading only the hypothesis reaches far above chance~\cite{Gururangan2018annotation,Poliak2018hypothesis}, while Niven and Kao~\cite{Niven2019probing} traced apparent argument comprehension to lexical cues that vanish under an adversarial split, with Geirhos et al.~\cite{Geirhos2020shortcut} giving the general account. Our surface-only arm is a hypothesis-only baseline transferred to this task, establishing how far a model gets without the capability the benchmark is taken to measure, except that it removes a property of the text rather than a span of it.

\textbf{Robustness and its controls.} Antoun et al.~\cite{Antoun2023towards} attack HC3-trained detectors with misspellings and homoglyphs, reporting a fall from 99.88 F1 to 33.57\% accuracy, and Huang et al.~\cite{Huang2024generated} report comparable degradation. Neither reports a label-free control, and our results suggest such collapses cannot be separated from a generic degenerate response to unreadable input without one, a methodological gap rather than a dispute about their measurements. Lai et al.~\cite{Lai2024adaptive} show ensembling preserves in-distribution accuracy while still degrading under shift, Zhang et al.~\cite{Zhang2024detecting} reframe detection as a two-sample test, and Borile and Abrate~\cite{Borile2025generalize} ablate roughly twenty feed-forward neurons for up to 6.9 points of out-of-distribution gain, evidence that some in-distribution performance rests on localised corpus-specific features.

\textbf{Fairness of the resulting decisions.} Liang et al.~\cite{Liang2023gpt} show several deployed detectors misclassify writing by non-native English speakers as machine-generated at high rates, which bears directly on ours. A detector whose separability comes substantially from orthographic regularity will concentrate its errors on writers whose orthographic habits differ from the corpus it was fitted to. We do not measure that effect here and note the connection because it is the most consequential reason to care which arm a benchmark's score comes from.

\textbf{Cross-corpus evaluation.} Alikhanov et al.~\cite{Alikhanov2026generated} merge HC3 and DAIGT V2 into one 124,195-sample corpus under a topic-held-out split, reporting 82.87\% to 88.86\% accuracy against the near-ceiling numbers usual in-distribution. Their design asks whether one model can serve both corpora, ours asks what each corpus is separable by, so the two sit alongside rather than replacing one another. Lu et al.~\cite{Lu2023large} show detectors can be evaded by prompting alone, a threat model orthogonal to the question here.